\subsection{等价}

设 \(\pmb{A}\) 和 \(\pmb{B}\) 为组装。组装

\[
(A \Longrightarrow B) \text { 且 } (B \Longrightarrow A)
\]

将被记为 \(A \Longleftrightarrow B\) .

CS7。设 \(A, B, T\) 为组装，且设 \(x\) 为一个字母。则组装 \((T|x)(A \Longleftrightarrow B)\) 与 \((T|x)A \Longleftrightarrow (T|x)B\) 相同。

这直接从CS5（§1，第2节）和CS6（第4节）得出。

CF10。如果A和B是C中的关系，则 \(A \Longleftrightarrow B\) 是T中的关系。

这直接从CF5（§1，第4节）和CF9（第4节）得出。

如果 \(A \leftrightarrow B\) 是 \(\mathfrak{C}\) 中的定理 ，我们将说 \(A\) 和 \(B\) 在 \(\mathfrak{C}\) 中等价；若 \(x\) 是一个不是 \(\mathfrak{C}\) 的常量的字母，并且如果 \(A\) 和 \(B\) 被视为 \(x\) 中的关系，那么 \(\mathfrak{C}\) 中的每个项满足其中一个也满足另一个。

由准则C20、C21（第4条）可知，为了证明 \(\mathfrak{C}\) 中形如 \(A \Longleftrightarrow B\) 的定理，必要且充分条件是能够证明 \(A \Longrightarrow B\) 和 \(B \Longrightarrow A\) 在 \(\mathfrak{C}\) 中。这通常这样来完成：先证明 \(B\)（在由 \(\mathfrak{C}\) 出发、通过添加公理 \(A\) 而得到的理论中），然后再证明 \(A\)（在由 \(\mathfrak{C}\) 出发、通过添加公理 \(B\) 而得到的理论中）。这些论述直接引出以下准则，其证明留给读者：

C22. 设A、B、C是T中的关系。如果 \(A \Longleftrightarrow B\) 是T中的定理，那么 \(B \Longleftrightarrow A\) 是T中的定理。如果 \(A \Longleftrightarrow B\) 并且 \(B \Longleftrightarrow C\) 是T中的定理，那么 \(A \Longleftrightarrow C\) 是T中的定理。

C23. 设A、B是C中的等价关系，且C是C中的关系。则以下各条是C中的定理：

\[
\begin{array}{c} \text {(非} A) \iff \text {(非} B); \quad (A \Rightarrow C) \iff (B \Rightarrow C); \\ (C \Rightarrow A) \iff (C \Rightarrow B); \end{array}
\]

\[
(A \text {   且   } C) \Longleftrightarrow (B \text {   且   } C); \quad (A \text {   或   } C) \Longleftrightarrow (B \text {   或   } C).
\]

C24. 设 \(A, B, C\) 是 \(\mathfrak{C}\) 中的关系。则以下各条是 \(\mathfrak{C}\) 中的定理：

（非非A） \(\Longleftrightarrow\) A；（A \(\Longrightarrow\) B） \(\Longleftrightarrow\) （（非B） \(\Longrightarrow\) （非A））；（A且A） \(\Longleftrightarrow\) A；（A且B） \(\Longleftrightarrow\) （B且A）；（A且（B且C）） \(\Longleftrightarrow\) （（A且B）且C）；（A或B） \(\Longleftrightarrow\) 非（（非A）且（非B））；（A或A） \(\Longleftrightarrow\) A；（A或B） \(\Longleftrightarrow\) （B或A）；（A或（B或C）） \(\Longleftrightarrow\) （（A或B）或C）；（A且（B或C）） \(\Longleftrightarrow\) （（A且B）或（A且C））；（A或（B且C）） \(\Longleftrightarrow\) （（A或B）且（A或C））；（A且（非B）） \(\Longleftrightarrow\) 非（A \(\Longrightarrow\) B）；（A或B） \(\Longleftrightarrow\) （（非A） \(\Longrightarrow\) B）。

C25. 如果A是C中的一个定理，且B是C中的一个关系，则

\[
(A \text {   且   } B) \Longleftrightarrow B
\]

是 \(\mathfrak{G}\) 中的定理 。如果``非 \(A\)''是 \(\mathfrak{G}\) 中的定理，那么 \((A\) 或 \(B) \Longleftrightarrow B\) 是 \(\mathfrak{G}\) 中的定理 .

¶ 原则上，从现在起直到本系列其余部分，准则C1至C25将不加引用地使用。

